\begin{proofbox}
	\textbf{\textcolor{CProof}{思路提示}}
	先处理过右焦点且斜率存在的情形。
	通过联立直线与椭圆方程，利用韦达定理和弦长公式，推导出焦点弦长公式。
	垂直直线再单独验证，左焦点由对称性得到同样结论。

	\medskip

	\textbf{\textcolor{CProof}{正式推导}}

	\begin{description}[style=nextline, leftmargin=2.8em, labelsep=0.5em, itemsep=0.55em]
		\item[\textbf{步骤一（设直线方程）：}]
		      先考虑过右焦点 $F(c,0)$ 且不垂直于 $x$ 轴的直线。
		      设直线方程为
		      \[
			      y=k(x-c).
		      \]
		      \par\textbf{理由：}
		      直线过点 $(c,0)$，且斜率为 $k$。

		\item[\textbf{步骤二（联立椭圆）：}]
		      代入椭圆
		      \[
			      \frac{x^2}{a^2}+\frac{y^2}{b^2}=1
		      \]
		      得
		      \[
			      \frac{x^2}{a^2}+\frac{k^2(x-c)^2}{b^2}=1.
		      \]
		      两边乘 $a^2b^2$，整理得
		      \[
			      (a^2k^2+b^2)x^2-2a^2k^2c\,x+(a^2k^2c^2-a^2b^2)=0.
		      \]
		      \par\textbf{理由：}
		      代入后去分母，并合并 $x$ 的同次项。

		\item[\textbf{步骤三（韦达定理）：}]
		      设两个交点的横坐标为 $x_1,x_2$，则
		      \[
			      x_1+x_2=\frac{2a^2k^2c}{a^2k^2+b^2},
			      \qquad
			      x_1x_2=\frac{a^2k^2c^2-a^2b^2}{a^2k^2+b^2}.
		      \]
		      \par\textbf{理由：}
		      二次方程 $Ax^2+Bx+C=0$ 的两根满足
		      \[
			      x_1+x_2=-\frac BA,\qquad x_1x_2=\frac CA.
		      \]

		\item[\textbf{步骤四（计算 $|x_1-x_2|^2$）：}]
		      \[
			      \begin{aligned}
				      (x_1-x_2)^2
				       & = (x_1+x_2)^2-4x_1x_2                               \\[0.3em]
				       & = \frac{4a^4k^4c^2}{(a^2k^2+b^2)^2}
				      -\frac{4(a^2k^2c^2-a^2b^2)}{a^2k^2+b^2}                \\[0.3em]
				       & = \frac{4a^4k^4c^2-4(a^2k^2c^2-a^2b^2)(a^2k^2+b^2)}
				      {(a^2k^2+b^2)^2}                                       \\[0.3em]
				       & = \frac{4a^2b^2(a^2k^2+b^2-k^2c^2)}
				      {(a^2k^2+b^2)^2}.
			      \end{aligned}
		      \]
		      又因为
		      \[
			      c^2=a^2-b^2,
		      \]
		      所以
		      \[
			      \begin{aligned}
				      a^2k^2+b^2-k^2c^2
				       & =a^2k^2+b^2-k^2(a^2-b^2) \\
				       & =b^2(1+k^2).
			      \end{aligned}
		      \]
		      因此
		      \[
			      (x_1-x_2)^2
			      =
			      \frac{4a^2b^4(1+k^2)}{(a^2k^2+b^2)^2}.
		      \]
		      \par\textbf{理由：}
		      代入韦达结果，通分后利用 $c^2=a^2-b^2$ 化简。

		\item[\textbf{步骤五（使用弦长公式）：}]
		      由于直线斜率为 $k$，两交点间距离满足
		      \[
			      |PQ|=\sqrt{1+k^2}\,|x_1-x_2|.
		      \]
		      由上式开方得
		      \[
			      |x_1-x_2|
			      =
			      \frac{2ab^2\sqrt{1+k^2}}{a^2k^2+b^2}.
		      \]
		      所以
		      \[
			      \begin{aligned}
				      |PQ|
				       & =\sqrt{1+k^2}\cdot
				      \frac{2ab^2\sqrt{1+k^2}}{a^2k^2+b^2} \\[0.3em]
				       & =\frac{2ab^2(1+k^2)}{a^2k^2+b^2}.
			      \end{aligned}
		      \]
		      \par\textbf{理由：}
		      弦长公式与韦达结果结合即可得到焦点弦长公式。
	\end{description}

	\medskip

	\textbf{\textcolor{CProof}{垂直情形验证}}
	当过右焦点的直线垂直于 $x$ 轴时，直线为
	\[
		x=c.
	\]
	代入椭圆方程得
	\[
		\frac{c^2}{a^2}+\frac{y^2}{b^2}=1.
	\]
	因为 $c^2=a^2-b^2$，所以
	\[
		\frac{y^2}{b^2}=1-\frac{c^2}{a^2}
		=1-\frac{a^2-b^2}{a^2}
		=\frac{b^2}{a^2}.
	\]
	故
	\[
		y=\pm\frac{b^2}{a}.
	\]
	于是通径长度为
	\[
		|PQ|=\frac{b^2}{a}-\left(-\frac{b^2}{a}\right)
		=\frac{2b^2}{a}.
	\]

	\medskip

	\textbf{\textcolor{CProof}{左焦点同理}}
	以上推导针对右焦点 $F(c,0)$。
	若直线过左焦点 $F(-c,0)$，可设直线为
	\[
		y=k(x+c).
	\]
	同样联立椭圆方程，或直接利用椭圆关于 $y$ 轴对称，可知所得弦长公式完全相同。

	\medskip

	\textbf{\textcolor{CProof}{角度形式推导}}
	若过焦点的直线与 $x$ 轴夹角为 $\theta$，且 $\cos\theta\ne0$，则
	\[
		k=\tan\theta.
	\]
	代入斜率公式：
	\[
		\begin{aligned}
			|PQ|
			 & =\frac{2ab^2(1+\tan^2\theta)}{a^2\tan^2\theta+b^2} \\[0.3em]
			 & =\frac{2ab^2\sec^2\theta}
			{a^2\frac{\sin^2\theta}{\cos^2\theta}+b^2}            \\[0.3em]
			 & =\frac{2ab^2}{a^2\sin^2\theta+b^2\cos^2\theta}.
		\end{aligned}
	\]
	当 $\theta=\frac{\pi}{2}$ 时，该式直接给出
	\[
		|PQ|=\frac{2ab^2}{a^2}=\frac{2b^2}{a},
	\]
	与通径公式一致。

	\medskip

	\textbf{\textcolor{CProof}{结论回扣}}
	因此，对于焦点在 $x$ 轴的椭圆
	\[
		\frac{x^2}{a^2}+\frac{y^2}{b^2}=1\qquad (a>b>0),
	\]
	过焦点且斜率为 $k$ 的非垂直直线对应的焦点弦长为
	\[
		|PQ|=\frac{2ab^2(1+k^2)}{a^2k^2+b^2};
	\]
	过焦点且垂直于 $x$ 轴的直线对应通径，弦长为
	\[
		|PQ|=\frac{2b^2}{a}.
	\]
\end{proofbox}